% Hochschild, coderived, and compact CY3 chiral scope.

\documentclass[11pt]{amsart}
\usepackage[localtheorems]{raeez-math-template}

\newtheorem{theorem}{Theorem}[section]
\newtheorem{proposition}[theorem]{Proposition}
\newtheorem{lemma}[theorem]{Lemma}
\newtheorem{conjecture}[theorem]{Conjecture}
\newtheorem{remark}[theorem]{Remark}

\newcommand{\HH}{\mathrm{HH}}
\newcommand{\HC}{\mathrm{HC}}
\newcommand{\HP}{\mathrm{HP}}
\newcommand{\THH}{\mathrm{THH}}
\newcommand{\TC}{\mathrm{TC}}
\newcommand{\ChirHoch}{\mathrm{ChirHoch}}
\newcommand{\kch}{\kappa_{\mathrm{ch}}}
\newcommand{\kcat}{\kappa_{\mathrm{cat}}}
\newcommand{\kBKM}{\kappa_{\mathrm{BKM}}}
\newcommand{\cO}{\mathcal{O}}
\newcommand{\cA}{\mathcal{A}}
\newcommand{\cC}{\mathcal{C}}
\newcommand{\Db}{D^{b}}
\newcommand{\Dco}{D^{\mathrm{co}}}
\newcommand{\Ran}{\mathrm{Ran}}
\newcommand{\MZV}{\mathrm{MZV}}
\newcommand{\Zchder}{Z^{\mathrm{der}}_{\mathrm{ch}}}
\DeclareMathOperator{\RHom}{RHom}

\begin{document}

\title{Hochschild, Coderived, and Compact CY3 Chiral Scope}
\author{}
\date{}

\begin{abstract}
The Hochschild and coderived constructions attached to a compact
Calabi-Yau threefold separate into five objects.  Hochschild homology
of $\Db(X)$ is an HKR invariant of the smooth proper category.
Chiral Hochschild cochains of $A=\Phi_3(\Db(X))$ compute the derived
chiral centre.  Positselski completion and Booth-Lazarev global Koszul
duality control curved bar-cobar inversion only after the algebraic
ambient is fixed.  For $d\geq3$ the native output of $\Phi_d$ is
$E_1$-chiral; the first $E_2$ object is the centre or double.
\end{abstract}

\maketitle

\tableofcontents

\section{The three Hochschild theories}

Let $X$ be a smooth proper complex variety.  The word Hochschild
denotes three different constructions.

\begin{enumerate}
\item $\HH_*(\Db(X))$ is Hochschild homology of the
$\mathbb C$-linear dg category of perfect complexes.  By
Hochschild-Kostant-Rosenberg,
\[
 \HH_k(\Db(X))
 \cong
 \bigoplus_{p-q=k} H^q(X,\Omega_X^p).
\]
The Connes operator produces the cyclic variants $\HC^-$ and $\HP$.

\item $\THH(\Db(X))$ is topological Hochschild homology.  Its output is
a cyclotomic spectrum.  After rationalization over a characteristic
zero base it shadows ordinary Hochschild theory; integrally it carries
trace-methods information invisible to $\HH_*$.

\item $\ChirHoch^\bullet(A)$ is chiral Hochschild cohomology of a
chiral algebra $A$ on a curve.  For $A=\Phi_d(\cC)$ it is a
factorisation cochain object on the Ran space.  It is not
$\HH_*(\cC)$ and it is not $\THH(\cC)$.
\end{enumerate}

\begin{proposition}[Typed Hochschild comparison]
\label{prop:typed-hochschild-comparison}
Let $\cC=\Db(X)$ and let $A_X=\Phi_3(\cC)$ be a stage-2
curve-specialised chiral algebra, when the construction exists.
There is no canonical equality
\[
 \HH_*(\cC)=\THH_*(\cC)=\ChirHoch^\bullet(A_X).
\]
Any comparison between $\HH^\bullet(\cC,\cC)$ and
$\ChirHoch^\bullet(A_X)$ requires the chosen $\Phi_3$ model,
the curve-specialisation datum, and compatibility with the
Gerstenhaber and Connes operations.
\end{proposition}

\begin{proof}
The three inputs and outputs live in different categories:
graded vector spaces with Connes operator for $\HH$, cyclotomic spectra
for $\THH$, and factorisation cochains for $\ChirHoch$.  A functor
$\Phi_3$ may transport selected Hochschild operations, but transport
is additional structure.  It is not a formal consequence of HKR and it
does not identify the $S^1$-action, conformal grading, and Ran-space
filtration.
\end{proof}

\section{HKR computation for \texorpdfstring{$\mathrm{K3}\times E$}{K3 x E}}

Let $S$ be a K3 surface and $E$ an elliptic curve.  The product
$X=S\times E$ is a compact Calabi-Yau threefold after choosing a
holomorphic volume form
\[
 \eta_X=\eta_S\wedge dz_E:
 \cO_X \xrightarrow{\sim} \omega_X .
\]
This orientation is part of the Calabi-Yau structure.  The equality
$\omega_X\simeq\cO_X$ alone does not choose the trace pairing.

\begin{lemma}[Product Hodge numbers]
\label{lem:k3e-hodge-product}
For smooth proper varieties,
\[
 h^{p,q}(Y\times Z)
 =
 \sum_{\substack{a+c=p\\ b+d=q}}
 h^{a,b}(Y)\,h^{c,d}(Z).
\]
For $S\times E$, the distribution of Hodge numbers by $p-q$ is
\[
\begin{array}{c|rrrrrrr}
k=p-q & -3&-2&-1&0&1&2&3\\ \hline
\sum_{p-q=k} h^{p,q}(S\times E) & 1&2&23&44&23&2&1 .
\end{array}
\]
\end{lemma}

\begin{proof}
The formula is Kunneth for Hodge cohomology.  The K3 distribution by
$p-q$ is $1,22,1$ in degrees $-2,0,2$.  The elliptic-curve
distribution is $1,2,1$ in degrees $-1,0,1$.  Their convolution gives
$1,2,23,44,23,2,1$.
\end{proof}

\begin{proposition}[Hochschild homology of \texorpdfstring{$S\times E$}{S x E}]
\label{prop:k3e-hh-hkr}
With the HKR convention
$\HH_k(\Db(X))=\bigoplus_{p-q=k}H^q(X,\Omega_X^p)$,
\[
\begin{array}{c|rrrrrrr}
k & -3&-2&-1&0&1&2&3\\ \hline
\dim \HH_k(\Db(S\times E)) & 1&2&23&44&23&2&1 .
\end{array}
\]
The total dimension is $96=\sum_{p,q}h^{p,q}(S\times E)$.
\end{proposition}

\begin{proof}
Apply HKR and Lemma~\ref{lem:k3e-hodge-product}.  The table is
palindromic by Serre duality after the orientation $\eta_X$ is fixed.
\end{proof}

\begin{remark}[Quintic check]
For a smooth quintic threefold $Q$, the nonzero Hodge numbers are
$h^{0,0}=h^{3,3}=h^{3,0}=h^{0,3}=1$,
$h^{1,1}=h^{2,2}=1$, and $h^{2,1}=h^{1,2}=101$.  Hence
\[
\dim\HH_k(\Db(Q))=(1,0,101,4,101,0,1)
\]
for $k=-3,\ldots,3$.  This is categorical Hochschild homology, not a
chiral Hochschild occupation statement.
\end{remark}

\begin{proposition}[Serre functor and Hochschild cochains]
\label{prop:serre-hh-cochains}
Let $X$ be a smooth proper oriented Calabi-Yau threefold.  Then
$S_{\Db(X)}=[3]$, and the chosen volume form gives
\[
 \HH^i(\Db(X))
 =
 \bigoplus_{p+q=i} H^p(X,\Lambda^q T_X)
 \cong
 \HH_{3-i}(\Db(X)).
\]
\end{proposition}

\begin{proof}
The Serre functor on $\Db(X)$ is
$(-)\otimes\omega_X[\dim X]$.  The orientation identifies it with
$[3]$.  Contraction with $\eta_X$ identifies
$\Lambda^qT_X$ with $\Omega_X^{3-q}$, and HKR gives the displayed
duality.
\end{proof}

\section{Bar, cobar, duality, and centre}

Let $A$ be the $E_1$ chiral algebra produced at stage 2 for a
Calabi-Yau input of dimension $d\geq3$.  Five associated objects must
remain distinct:
\[
\begin{array}{cl}
A & \text{the original }E_1\text{ chiral algebra},\\
B(A)=T^c(s^{-1}\bar A) & \text{the ordered bar coalgebra},\\
A^i=H^\bullet(B(A)) & \text{the bar-cohomology coalgebra},\\
A^!=D_{\Ran}(A^i) & \text{the Verdier-dual Koszul algebra},\\
\Zchder(A)=C^\bullet_{\mathrm{ch}}(A,A)
  \simeq \RHom(\Omega B(A),A) & \text{the derived chiral centre}.
\end{array}
\]

\begin{proposition}[No collapse of the five objects]
\label{prop:five-object-separation}
The equivalence $\Omega B(A)\simeq A$, when valid, is bar-cobar
inversion.  It does not identify $B(A)$ with $A^i$, $A^!$, or
$\Zchder(A)$.
\end{proposition}

\begin{proof}
The four functors after $A$ are different: bar construction, bar
cohomology, Verdier duality on the Ran side, and chiral Hochschild
cochains.  Applying one of them does not apply the others.  The object
$\Omega B(A)$ is a cobar recovery of $A$; the centre
$\Zchder(A)$ is an endomorphism object of $A$ as a bimodule.
\end{proof}

\section{Coderived completion}

Curvature breaks the standard derived category.  In a curved algebra
or coalgebra the differential satisfies $d^2=[h,-]$ or its coalgebraic
dual.  Ordinary cohomology is then the wrong invariant.  Positselski's
coderived and contraderived categories replace ordinary quasi-isomorphisms
by maps whose cones are coacyclic or contraacyclic after the appropriate
bar-cobar test.

\begin{theorem}[Booth-Lazarev, associative scope]
\label{thm:booth-lazarev-scope}
Booth and Lazarev construct model structures on curved coalgebras and
curved algebras and prove that the extended bar-cobar adjunction is a
Quillen equivalence.  In the conilpotent coalgebra and uncurved dg
algebra subcase this recovers the ordinary Positselski and
Keller-Lefevre Koszul duality.
\end{theorem}

\begin{proof}
This is the associative theorem of Booth-Lazarev, \emph{Global Koszul
duality}, arXiv:2304.08409, revised version v4 of 5 January 2026.
The theorem is stated for curved algebraic objects, not for
factorisation coalgebras on $\Ran(C)$.
\end{proof}

\begin{conjecture}[Chiral Booth-Lazarev lift]
\label{conj:chiral-booth-lazarev}
Let $A$ be a curved $E_1$ chiral algebra on a smooth curve $C$.  A
chiral Booth-Lazarev theorem would give a Quillen equivalence
\[
 \Dco_{\mathrm{ch}}(B_{\mathrm{curv}}(A))
 \simeq
 D^{\mathrm{fact}}(\Omega_{\mathrm{curv}}B_{\mathrm{curv}}(A))
\]
on curved factorisation coalgebras over $\Ran(C)$.
This is not a consequence of Theorem~\ref{thm:booth-lazarev-scope}.
\end{conjecture}

\begin{proof}[Proof obligation]
The associative theorem supplies the algebraic model.  The chiral lift
requires three additional verifications.
\begin{enumerate}
\item A cofibrantly generated model structure on curved factorisation
coalgebras over $\Ran(C)$, compatible with the factorisation tensor
product and the Ran topology.
\item A pro-conilpotent completion filtration respected by the chiral
bar differential, so that $B(A)$ is replaced by
$\widehat B(A)=\varprojlim_w B(A)/F^w$ whenever ordinary conilpotence
fails.
\item Boundary descent for the curvature family
$m_0^{g,n}\in\Gamma(\overline{\mathcal M}_{g,n},
\lambda_g\otimes H^2(A))$ under clutching maps.  Booth-Lazarev handle
a curved algebraic object; the moduli-indexed sewing family is extra
data.
\end{enumerate}
Until these three statements are proved, chiral Booth-Lazarev remains a
conjecture with a precise associative shadow.
\end{proof}

\section{Native operadic level at \texorpdfstring{$d\geq3$}{d >= 3}}

\begin{proposition}[E1 output, E2 centre]
\label{prop:e1-output-e2-centre}
For a Calabi-Yau input of dimension $d\geq3$, the stage-2 output
$A=\Phi_d(\cC)$ is native $E_1$ chiral.  The first $E_2$ object
attached to $A$ is the centre:
\[
 \mathcal Z(\mathrm{Rep}^{E_1}(A))
 \simeq
 \mathrm{Rep}^{E_2}(\Zchder(A)),
\]
under the usual Morita-dualisability and presentability hypotheses.
\end{proposition}

\begin{proof}
The $E_2$ structure is supplied by the higher Deligne theorem on
Hochschild cochains and by the Drinfeld-centre comparison for the
$E_1$ representation category.  Both constructions take $A$ as input.
Neither changes the native operadic level of $A$.
\end{proof}

\section{Compact CY3 theorem obligations}

A compact Calabi-Yau threefold does not automatically produce a proved
chiral theorem.  The following data are part of the theorem, not
decoration.

\begin{proposition}[Minimum data for compact CY3 chiral output]
\label{prop:compact-cy3-obligations}
Let $X$ be a compact Calabi-Yau threefold.  A proof that
$A_X=\Phi_3(\Db(X))$ exists as a quantum $E_1$ chiral algebra requires:
\begin{enumerate}
\item an orientation $\eta_X:\cO_X\xrightarrow{\sim}\omega_X$, or
equivalently a Calabi-Yau trace class making the Serre functor $[3]$;
\item stage-2 specialisation data $(\Sigma_2,C)$ and descent of the
factorisation algebra to the chosen curve;
\item a solution of the quantum master equation, including anomaly
cancellation and a renormalisation scheme;
\item Positselski or Booth-Lazarev completion whenever the bar
coalgebra is not conilpotent;
\item Ran-space and boundary-clutching descent for any genus-indexed
curvature terms.
\end{enumerate}
Without these hypotheses, only the categorical HKR and Serre-duality
statements of Sections~1 and~2 are proved.
\end{proposition}

\begin{proof}
Items (1) and (2) define the oriented categorical and chiral input.
Item (3) is the obstruction-theoretic condition for quantisation of the
factorisation algebra.  Item (4) identifies the correct derived ambient
for curved bar-cobar inversion.  Item (5) is needed because compact CY3
curvature data appears as a sewing-compatible family, not as a single
scalar curvature.
\end{proof}

\section{Automorphic constants used here}

The Hochschild and coderived arguments use only the following
automorphic normalisations.
\[
 \kcat(S\times E)
 =
 \chi(\cO_S)\chi(\cO_E)
 =
 2\cdot0
 =
 0.
\]
For the Gritsenko form $\Delta_5$, the constant coefficient
normalisation is
\[
 c_1(0)=10,\qquad
 \kBKM(\Delta_5)=c_1(0)/2=5.
\]
The Mukai-Heisenberg invariant
\[
 K^{\kch}=8
\]
belongs to the K3 Mukai-Heisenberg sector.  It is not
$\kBKM(\Delta_5)$ and it is not a general identity
$\kch+\kch^!=8$ for compact CY3 outputs.

\begin{remark}[Constants not proved here]
The five-archetype ceiling
$\kch+\kch^!\in\{0,8,13,250/3,98/3\}$ and the Monster row require the
separate Theorem C calculation.  The present Hochschild-coderived
argument supplies no proof of the values $250/3$, $214/3$, $98/3$, or
$74/3$.
\end{remark}

\section{Single-valued periods}

The only period statement needed for the Hochschild-coderived scope is
the following conditional restriction.

\begin{proposition}[Single-valued restriction]
\label{prop:sv-restriction}
If a chiral Hochschild period comparison factors through Brown's
single-valued motivic period map, then
\[
 \zeta(2)\mapsto0,\qquad
 \zeta(2r+1)\mapsto 2\zeta(2r+1).
\]
Consequently a weight-2 coefficient does not survive on the
single-valued target.
\end{proposition}

\begin{proof}
This is Brown's single-valued period map.  No depth-two or depth-three
multiple-zeta relation is used in this note.  Such a relation becomes
load-bearing only after an explicit configuration-space integral is
computed and matched to the chiral Hochschild complex.
\end{proof}

\section{Analytic bar Laplacian}

\begin{lemma}[A sufficient self-adjointness criterion]
\label{lem:bar-laplacian-sufficient}
Let
$\mathcal H=\widehat{\bigoplus}_{n\geq0}\mathcal H_n$ be a Hilbert
direct sum with each $\mathcal H_n$ finite-dimensional.  Let
$D=\bigoplus_n\mathcal H_n$ be the finite-support domain.  If a
symmetric non-negative operator $\Delta$ preserves each $\mathcal H_n$
and the restriction $\Delta|_{\mathcal H_n}$ is self-adjoint, then
$\Delta|_D$ is essentially self-adjoint.
\end{lemma}

\begin{proof}
The closure is the Hilbert direct sum of the finite-dimensional
self-adjoint matrices $\Delta|_{\mathcal H_n}$.  The finite-support
domain is a core for this direct sum.
\end{proof}

\begin{remark}[What remains to prove for a CY3 bar Laplacian]
Finite-dimensional conformal-weight pieces do not by themselves prove
essential self-adjointness.  One must construct the Hilbert completion,
prove that the bar differential and its adjoint preserve the block
decomposition, and prove symmetry and non-negativity of
$d_Bd_B^*+d_B^*d_B$ on the chosen domain.  These analytic choices are
independent of the algebraic coderived equivalence.
\end{remark}

\section*{References}

\begin{itemize}
\item Beilinson, A. and Drinfeld, V.,
\emph{Chiral Algebras}, AMS Colloquium Publications 51, 2004.
\item Booth, M. and Lazarev, A.,
\emph{Global Koszul duality}, arXiv:2304.08409, v4, 2026.
\item Brown, F.,
\emph{Single-valued periods and multiple zeta values},
Forum of Mathematics, Sigma 2 (2014), e25.
\item Caldararu, A. and Willerton, S.,
\emph{The Mukai pairing, I: a categorical approach},
New York Journal of Mathematics 16 (2010), 61-98.
\item Francis, J.,
\emph{The tangent complex and Hochschild cohomology of $E_n$-rings},
Compositio Mathematica 149 (2013), 430-480.
\item Gritsenko, V. and Nikulin, V.,
\emph{Siegel automorphic form corrections of some Lorentzian Kac-Moody
Lie algebras}, American Journal of Mathematics 120 (1998), 1-54.
\item Malikov, F., Schechtman, V. and Vaintrob, A.,
\emph{Chiral de Rham complex}, Communications in Mathematical Physics
204 (1999), 439-473.
\item Mukai, S.,
\emph{Symplectic structure of the moduli space of sheaves on an
abelian or K3 surface}, Inventiones Mathematicae 77 (1984), 101-116.
\item Pantev, T., Toen, B., Vaquie, M. and Vezzosi, G.,
\emph{Shifted symplectic structures}, Publications Mathematiques de
l'IHES 117 (2013), 271-328.
\item Positselski, L.,
\emph{Two kinds of derived categories, Koszul duality, and comodule
contramodule correspondence}, Memoirs of the AMS 212 (2011), no. 996.
\end{itemize}

\end{document}
